% !PRINTED 67
\weberpage{67}
That in the case of an even $\beta$ each of the two factors vanishes is seen by setting, in (15),
\[
n' = 1 + 2^{\lambda-1},
\]
and consequently
\[
r^{n'} = -r, \qquad \nu_1' = 0, \qquad \nu' = 2^{\lambda-3}
\]
there. Then this formula gives, since $n$ is odd:
\[
\bigl((-1)^{\beta_1}, \Theta^\beta, r\bigr) = -\Theta^{\beta 2^{\lambda-3}} \bigl((-1)^{\beta_1}, \Theta^\beta, r\bigr).
\]
For even $\beta$, however, $\Theta^{\beta 2^{\lambda-3}} = +1$, and consequently:
\weq{18}{\bigl((-1)^{\beta_1}, \Theta^\beta, r\bigr) = 0;}
hence:

2.~\gesperrt{If the degree of the root of unity $r$ is a power of 2 and greater than 4, then the resolvent $\bigl((-1)^{\beta_1}, \Theta^\beta, r\bigr)$ vanishes if and only if $\beta$ is even.}

In the case $m=4$, and so $r=i$, one has only the two resolvents $i+i^3, i-i^3$, which we can combine under the designation $((-1)^\beta, i)$, $\beta=0,1$, and $((-1)^\beta, i)$ is $=0$ if and only if $\beta=0$.

\begin{center}
\S.\,18.\\
\gesperrt{Cyclotomic Fields.}
\end{center}

The theory of \gesperrt{Abelian} groups opens to us a deeper insight into the theory of the roots of unity and of the algebraic numbers that spring from them.

Now let $m$ be any positive integer, decomposed into its prime factors:
\weq{1}{m = 2^\lambda q_1^{\varkappa_1} q_2^{\varkappa_2} \cdots,}
and let $r$ be a primitive $m$th root of unity. The case $\lambda=1$ we can exclude once and for all from our consideration; for if $m$ is odd and $r$ runs through the primitive $m$th roots of unity, no two opposite ones occur among them, and $-r$ runs through the primitive $2m$th roots of unity.

$r$ is the root of an irreducible \gesperrt{Abelian} equation with integral coefficients, of degree
\weq{2}{\nu = \varphi(m),}
as we have proved in \S.\,134 of the first volume.
